\vsssub
\subsubsection{~多网格} \label{sec:mgrids}
\vsssub

到目前为止，只考虑了单个波浪模式网格。为使单个程序能够运行多个模式网格，
需要设计一种数据结构，在其中同时保留所有不同模式网格以及所有模式的内部
工作数组，并提供一个简单机制来选择实际要操作的波浪模式网格。为实现这一点，
按如下方式使用若干 FORTRAN 90 特性 \citep[例如][]{bk:MR99}：

\begin{list}{}{\rightmargin 8mm \leftmargin 10mm \labelsep 2mm}

\item [1)] 使用 {\F type} 声明，在模式代码中定义一个或多个数据结构，
           其中包含模式设置和相关工作数组。

\item [2)] 构造这些数据结构的数组，数组中的每个元素定义一个独立模式网格。

\item [3)] 把描述模式的基本参数（例如网格点数 {\F nx} 和 {\F ny}）重新定义
           为指针，并让它们指向相应数据结构中的相应元素，从而生成即时别名。

\end{list}

\noindent
通过这种方式，可以定义多模式数据结构，同时在模式子程序内部保持描述模式的
所有原始变量布局不变。图~\ref{fig:struc_1} 和图~\ref{fig:struc_2} 用实际
源代码中的示例说明了这种结构及其使用。注意，结构内部类似 {\F zb} 的指针
数组通过 \command{\F allocate grids(imod)\%zb(nsea)} 分配内存。
\begin{figure}
\begin{center} \begin{tabular}{|c|} \hline 
\begin{minipage}[t]{4.5in}
\begin{verbatim}

!/     
!/ Data structures
!/                              
      TYPE GRID
        INTEGER               :: NX, NY, NSEA
        REAL, POINTER         :: ZB(:)
      END TYPE GRID
!/     
!/ Data storage
!/     
      TYPE(GRID), TARGET, ALLOCATABLE :: GRIDS(:)
!/    
!/ Data aliasses
!/ 
      INTEGER, POINTER        :: NX, NY, NSEA
      REAL, POINTER           :: ZB(:):
!/

\end{verbatim}
\end{minipage} \\ \hline
\end{tabular} \end{center}

\caption{{\file w3gdatmd.ftn} 中用于定义波浪模式多个空间网格的数据结构声明
示例。为简化起见，该示例只考虑网格维度 {\F nx}、{\F ny} 和 {\F nsea}，
以及底深数组 {\F zb}。}
\label{fig:struc_1}

\botline
\end{figure}

执行该语句后，别名指针 {\F zb} 还需要再次指向结构中的相应元素，才能让
该别名正确指向新分配的空间。因此，为该结构分配数组的子程序 {\F w3dimx}
在末尾包含一个对 {\F w3setx} 的调用；后者会为所选网格设置所有指针别名。
其他结构中设置数组大小的子程序同样如此。

\begin{figure}
\begin{center} \begin{tabular}{|c|} \hline
\begin{minipage}[t]{4.5in}
\begin{verbatim}

!
      NX     => GRIDS(IMOD)%NX
      NY     => GRIDS(IMOD)%NY
      NSEA   => GRIDS(IMOD)%NSEA
!
      ZB     => GRIDS(IMOD)%ZB
!

\end{verbatim}
\end{minipage} \\ \hline
\end{tabular} \end{center}

\caption{用于为模式编号 {\F imod} 激活图~\ref{fig:struc_1} 中指针别名的
源代码示例。}
\label{fig:struc_2}

\botline
\end{figure}

